Enterprise document corpora routinely store facts about the same entity in
different document types, written by different teams at different times. We ask
whether a knowledge graph keyed on stable entity identifiers retrieves such facts
better than similarity search over embeddings, and we measure the answer rather
than assert it. Our corpus is one snapshot of SEC EDGAR \citep{sec2024edgar}:
167{,}050 filings from 10{,}511 companies, 508{,}714 embedded chunks of
risk-factor text, and a graph of 430{,}257 nodes and 561{,}144 edges carrying an
accession number on every edge. Three architectures --- text-to-SQL over typed
columns, dense-vector RAG with HNSW cosine search
\citep{lewis2020rag,malkov2018hnsw}, and a two-stage GraphRAG design --- share one
model, temperature and token budget, and are scored by a programmatic judge over
20 questions in four categories, each asked three times, for 180 scored attempts.
GraphRAG solved 14 of 20 questions on all three attempts to text-to-SQL's 9 and
vector RAG's 3. Two results contradict the pre-registered expectation. Vector RAG
lost the single-document category it was expected to win, passing 3 of 15 attempts
where GraphRAG passed 15 of 15. Text-to-SQL beat GraphRAG on cross-document entity
resolution, 14 of 15 attempts against 9 of 15, in a category written for the
graph. GraphRAG also lost structured aggregation as predicted, 6 of 15 against
15 of 15, at 24 times the token cost. Only 2 of the 180 attempts asserted
something verifiably false, both text-to-SQL conflating distinct people who share
a surname. We conclude that structural retrieval earns its ingestion and token
cost on questions turning on identity or on joins across documents, and not
otherwise; we claim no advantage beyond this corpus, this model and this question
set.
